\begin{table}[htbp]
\centering
\caption{\\ Stock market responses to other legal and political events}
\label{tab:other_events}
\footnotesize
\setlength{\tabcolsep}{4pt}
\medskip
\begin{adjustbox}{max width=\textwidth}
\begin{tabular}{lccccccc}
\toprule
 & & & & \multicolumn{2}{c}{CAR (equal-wt.)} & & \\
\cmidrule(lr){5-6}
Event & Dates & Days & Market & Gold & No gold & Diff. & $t$ \\
\midrule
First gold-clause hearing (Irving Trust) & May 22--24, 1933 & 3 & +5.5\% & +2.8\% & +2.1\% & +0.7\% & 0.8 \\
\textit{In re Missouri Pacific} ruling & June 20--22, 1934 & 3 & -3.8\% & -2.1\% & -1.2\% & -0.9\% & -1.0 \\
\textit{Norman} affirmed (N.Y. Ct.\ App.) & July 3--6, 1934 & 3 & +2.2\% & -0.2\% & -0.5\% & +0.3\% & 0.4 \\
Cert.\ granted, \textit{Norman} & Oct.\ 8--10, 1934 & 3 & +0.7\% & +1.1\% & +0.5\% & +0.6\% & 0.6 \\
Cert.\ granted \& midterm election & Nov.\ 5--8, 1934 & 3 & +2.1\% & +5.1\% & +2.0\% & +3.2\% & 3.6 \\
\midrule
Joint Resolution (reference) & May 26--June 6, 1933 & 9 & +12.1\% & +15.9\% & +12.7\% & +3.2\% & 2.1 \\
\bottomrule
\end{tabular}
\end{adjustbox}
\medskip
\begin{minipage}{0.95\textwidth}\footnotesize \textit{Notes.} This table reports stock market responses to the intermediate legal and political events of the gold clause litigation, using the same portfolios, estimation window, and CAR methodology as Table~\ref{tab:event_study} in the main text. ``Diff.''\ is the difference between the cumulative abnormal returns of the gold-exposed and non-exposed equal-weighted portfolios. The $t$-statistic scales this difference by $\hat\sigma\sqrt{h}$, where $h$ is the number of trading days in the window and $\hat\sigma = 0.50$ percentage points is the standard deviation of the daily abnormal-return differential over calm trading days in 1934 (excluding the event windows in this table). The Irving Trust hearing (May~22--24, 1933) was the first court hearing on the enforceability of gold clauses. \textit{In re Missouri Pacific} (E.D.~Mo., June~20, 1934) was the first federal ruling upholding the constitutionality of the Joint Resolution; the New York Court of Appeals affirmed \textit{Norman v.\ Baltimore \& Ohio R.~Co.} on July~3, 1934. Certiorari was granted in \textit{Norman} on October~8, 1934 and in \textit{United States v.\ Bankers Trust Co.} on November~5, 1934; the New York Stock Exchange was closed on November~6, 1934 for the midterm election, so the November window mixes the certiorari grant with the election result. The Joint Resolution window from Table~\ref{tab:event_study} is repeated for reference.\end{minipage}
\end{table}
